% ===== PRÉAMBULE CANONIQUE — à coller VERBATIM en tête de CHAQUE .tex =====
\documentclass[11pt,a4paper]{article}
\usepackage[utf8]{inputenc}\usepackage[T1]{fontenc}\usepackage{lmodern}
\usepackage[french]{babel}
\usepackage{amsmath,amssymb,mathtools}\usepackage{icomma}
\usepackage[version=4]{mhchem}          % \ce{...} : équations & formules
\usepackage{siunitx}\sisetup{locale=FR} % \qty{}{}, \si{}, \ang{}
\usepackage{chemfig}                    % \chemfig{...} : structures développées
\usepackage{xcolor}\usepackage{colortbl}
\usepackage{tikz}\usepackage{pgfplots}\pgfplotsset{compat=1.18}
\usepackage[most]{tcolorbox}\usepackage{enumitem}\usepackage{array,booktabs}
\usepackage[a4paper,margin=2.2cm]{geometry}
\usetikzlibrary{arrows.meta,calc,angles,quotes,positioning,patterns,backgrounds,shapes.geometric}
\definecolor{primary}{RGB}{222,49,99}\definecolor{primarydark}{RGB}{150,22,55}
\definecolor{defcol}{RGB}{0,114,178}\definecolor{methcol}{RGB}{52,92,148}
\definecolor{excol}{RGB}{0,135,90}\definecolor{attncol}{RGB}{200,40,55}
\tcbset{boxbase/.style={enhanced,breakable,boxrule=0.9pt,arc=2.5pt,
  left=3mm,right=3mm,top=2.3mm,bottom=2.3mm,fonttitle=\bfseries,coltitle=white}}
% --- boîtes TUTORIEL (noms EXACTS, ne pas renommer) ---
\newtcolorbox{cle}[1][]{boxbase,colback=defcol!6,colframe=defcol,
  title={\textsc{À retenir}\ifx\relax#1\relax\relax\else\textmd{ : \itshape #1}\fi}}
\newtcolorbox{methode}[1][]{boxbase,colback=methcol!7,colframe=methcol,sharp corners,
  title={\textsc{Méthode}\ifx\relax#1\relax\relax\else\textmd{ --- \itshape #1}\fi}}
\newtcolorbox{exemplebox}[1][]{boxbase,colback=excol!5,colframe=excol,
  title={\textsc{Exemple}\ifx\relax#1\relax\relax\else\textmd{ : \itshape #1}\fi}}
\newtcolorbox{piege}[1][]{boxbase,colback=attncol!6,colframe=attncol,
  title={\textsc{Piège}\ifx\relax#1\relax\relax\else\textmd{ : \itshape #1}\fi}}
\newtcolorbox{remarque}[1][]{boxbase,colback=black!4,colframe=black!45,
  title={\textsc{Remarque}\ifx\relax#1\relax\relax\else\textmd{ : \itshape #1}\fi}}
% --- boîtes EXERCICE / CORRIGÉ (noms EXACTS, auto-comptées) ---
\newtcolorbox[auto counter]{exo}[1][]{boxbase,colback=primary!4,colframe=primary,
  title={\textsc{Exercice}~\thetcbcounter\ifx\relax#1\relax\relax\else\textmd{ --- \itshape #1}\fi}}
\newtcolorbox[auto counter]{corrige}[1][]{boxbase,colback=excol!4,colframe=excol,
  title={\textsc{Corrigé}~\thetcbcounter\ifx\relax#1\relax\relax\else\textmd{ --- \itshape #1}\fi}}
\newcommand{\R}{\mathbb{R}}\newcommand{\N}{\mathbb{N}}\newcommand{\Z}{\mathbb{Z}}\newcommand{\ds}{\displaystyle}
% (un \usepackage{fancyhdr}+en-tête est possible ; il est ignoré côté web)
% =========================================================================
\begin{document}
% THEME_TITLE: L'équation d'état pV=nRT et le volume molaire aux CNTP

Une seule équation relie les quatre variables d'un gaz parfait : l'\textbf{équation d'état}. Ce thème
entraîne à l'utiliser dans les deux sens (trouver $n$, $p$, $V$ ou $T$) et à exploiter le raccourci du
\textbf{volume molaire} aux conditions normales.

\begin{cle}[Équation d'état des gaz parfaits]
\[ \boxed{\,p\,V = n\,R\,T\,} \]
$p$ pression, $V$ volume, $n$ quantité de matière (\si{\mole}), $T$ température absolue (\si{\kelvin}),
$R$ constante des gaz parfaits. On en isole n'importe quelle grandeur :
\[ n=\frac{pV}{RT},\qquad p=\frac{nRT}{V},\qquad V=\frac{nRT}{p},\qquad T=\frac{pV}{nR}. \]
\end{cle}

\begin{cle}[Choisir la bonne valeur de $R$]
\begin{center}\renewcommand{\arraystretch}{1.4}
\begin{tabular}{c c c c}
\toprule
\rowcolor{primary!12}
\textbf{Valeur de $R$} & \textbf{$p$} & \textbf{$V$} & \textbf{$T$}\\
\midrule
$\qty{0.082}{atm.L.mol^{-1}.K^{-1}}$ & \si{atm} & \si{\litre} & \si{\kelvin}\\
$\qty{8.314}{\joule\per\mole\per\kelvin}$ & \si{\pascal} & \si{\meter\cubed} & \si{\kelvin}\\
\bottomrule
\end{tabular}
\end{center}
On \emph{ne panache jamais} : on choisit une ligne et on convertit toutes les données pour lui
correspondre. En chimie, $R=\num{0,082}$ (couple \si{atm}/\si{\litre}) est le plus commode.
\end{cle}

\begin{cle}[Conditions normales (CNTP) et volume molaire]
Aux \textbf{CNTP} (= TPN = PTN) : $p=\qty{1}{atm}$ et $T=\qty{273}{\kelvin}$. Alors \textbf{une mole de
n'importe quel gaz parfait} occupe le \textbf{volume molaire}
\[ \boxed{\,V_{\mathrm{m}}=\frac{RT}{p}=\qty{22.4}{\litre\per\mole}\,}
   \qquad\Longrightarrow\qquad
   n=\frac{V}{\qty{22.4}{\litre\per\mole}},\quad V=n\times\qty{22.4}{\litre\per\mole}. \]
\end{cle}

\begin{center}
\begin{tikzpicture}[box/.style={draw,thick,rounded corners=2pt,align=center,inner sep=5pt,font=\small},
   >=Stealth,node distance=8mm]
  \node[box,fill=defcol!10] (m) {masse $m$ (\si{\gram})};
  \node[box,fill=primary!10,right=22mm of m] (n) {quantité $n$ (\si{\mole})};
  \node[box,fill=excol!12,right=22mm of n] (v) {volume $V$ (\si{\litre}, CNTP)};
  \draw[->] (m) -- node[above,font=\scriptsize]{$\div\,M$} (n);
  \draw[->] (n) -- node[above,font=\scriptsize]{$\times\,22{,}4$} (v);
  \draw[->] (v) to[bend left=30] node[below,font=\scriptsize]{$\div\,22{,}4$} (n);
  \draw[->] (n) to[bend left=30] node[below,font=\scriptsize]{$\times\,M$} (m);
\end{tikzpicture}
\end{center}

\begin{methode}[Résoudre un problème avec $pV=nRT$]
\begin{enumerate}[leftmargin=1.6em,topsep=2pt,itemsep=1pt]
  \item Convertir : $T$ en \si{\kelvin} ; $p$, $V$ cohérents avec le $R$ choisi.
  \item Isoler l'inconnue littéralement.
  \item Application numérique, puis vérifier l'ordre de grandeur.
  \item \textbf{Cas CNTP} ($\qty{1}{atm}$, \qty{273}{\kelvin}) : passer par $V_{\mathrm{m}}=\qty{22.4}{\litre\per\mole}$
        est plus rapide.
\end{enumerate}
\end{methode}

\begin{exemplebox}[deux calculs types]
$n=\qty{0.50}{\mole}$ de \ce{O2} dans $V=\qty{5.0}{\litre}$ à \qty{27}{\degreeCelsius} :
$p=\dfrac{nRT}{V}=\dfrac{0{,}50\times0{,}082\times300}{5{,}0}=\qty{2.46}{atm}$. \\[2pt]
$V=\qty{44.8}{\litre}$ de gaz aux CNTP : $n=\dfrac{44{,}8}{22{,}4}=\qty{2.0}{\mole}$.
\end{exemplebox}

\begin{piege}[le volume molaire n'est valable qu'aux CNTP]
$\qty{22.4}{\litre\per\mole}$ suppose $p=\qty{1}{atm}$ \emph{et} $T=\qty{273}{\kelvin}$. Hors de ces
conditions (autre pression, autre température), il faut repasser par $pV=nRT$. Et toujours : $T$ en
kelvins !
\end{piege}

\end{document}
